\chapter{Yêu Thương Không Điều Kiện (Unconditional Love)}
\begin{center}
  \textit{\color{truyengold}``Tình yêu thương chân thật không đặt ra điều kiện --- nó chỉ cho đi và cho đi mãi.''}\\[4pt]
  \textit{(True love sets no conditions --- it simply gives and keeps giving.)}\\[4pt]
  \small\color{truyenblue}--- Thiền học Phật giáo ---
\end{center}
\ngancach

\section{Người Cha Đứng Chờ Ở Cuối Đường}
\begin{truyen}{Người Cha Đứng Chờ Ở Cuối Đường}{Kinh Thánh --- Luca 15:11--32}
\chuhoa{T}{rong} câu chuyện đứa con hoang đàng, người con thứ đòi chia gia sản rồi bỏ đi ăn chơi phá sạch tất cả. Khi cùng đường, anh ta quyết định quay về.\\[4pt]
\textit{(In the parable of the prodigal son, the younger son demanded his inheritance and left to squander it all in revelry. When he hit rock bottom, he decided to return home.)}

``Còn cách nhà xa, cha anh đã trông thấy --- và chạy ra ôm anh hôn.''\\[4pt]
\textit{(`While he was still a great way off, his father saw him --- and ran to embrace and kiss him.')}

Người cha không nói ``Tao đã bảo mà'' --- không đặt điều kiện --- không trách móc. Ông lập tức tổ chức tiệc mừng.\\[4pt]
\textit{(The father did not say `I told you so' --- set no conditions --- made no reproach. He immediately held a celebration.)}

Đây là hình ảnh của tình yêu thương trọn vẹn: luôn chờ đợi, luôn sẵn sàng, luôn mừng vui khi người lầm đường biết quay về.\\[4pt]
\textit{(This is the image of complete love: always waiting, always ready, always rejoicing when the one who went astray finds their way back.)}

\end{truyen}
\begin{baihoc}
  Yêu thương thật sự không tính toán quá khứ --- nó chỉ nhìn vào hiện tại và tương lai với trái tim rộng mở.\\[4pt]
  \textit{(True love does not calculate the past --- it only looks at the present and future with an open heart.)}
\end{baihoc}

\section{Bà Mẹ Chèo Đò Nuôi Con Ăn Học}
\begin{truyen}{Bà Mẹ Chèo Đò Nuôi Con Ăn Học}{Câu chuyện Việt Nam --- Thế kỷ 20}
\chuhoa{T}{rong} một làng chài ven sông, người mẹ góa bụa mỗi ngày thức dậy từ 3 giờ sáng, chèo đò chở khách đến chiều tối, để đủ tiền cho con đi học ở thành phố.\\[4pt]
\textit{(In a fishing village by the river, a widowed mother woke each day at 3 a.m., ferried passengers from dawn until evening, to earn enough for her son to study in the city.)}

Người con không biết mẹ vất vả đến thế --- mẹ không nói, chỉ gửi tiền đều đặn và viết: ``Con cứ học cho tốt là mẹ vui rồi.''\\[4pt]
\textit{(The son did not know how hard his mother worked --- she never said, only sent money regularly and wrote: `Just study well and mother is happy.')}

Ngày con đỗ đại học, người ta thấy bà đứng một mình ở bến đò, nhìn về phía thành phố, mỉm cười --- đôi bàn tay chai sần vì mái chèo, đôi mắt ngập tràn nước mắt hạnh phúc.\\[4pt]
\textit{(The day her son passed the university entrance exam, she was seen standing alone at the ferry dock, gazing toward the city, smiling --- her hands calloused from the oar, her eyes overflowing with tears of happiness.)}

\end{truyen}
\begin{baihoc}
  Tình mẹ không ồn ào --- nó lặng lẽ như bàn tay chèo đò mỗi sáng sớm, kiên trì như dòng sông không bao giờ ngừng chảy.\\[4pt]
  \textit{(A mother's love is not loud --- it is quiet as oar-strokes every early morning, persistent as a river that never stops flowing.)}
\end{baihoc}

\section{Janusz Korczak Và 192 Đứa Trẻ}
\begin{truyen}{Janusz Korczak Và 192 Đứa Trẻ}{Janusz Korczak --- Ba Lan, 1942}
\chuhoa{J}{anusz} Korczak là bác sĩ và nhà giáo dục nổi tiếng người Do Thái ở Warsaw; ông có thể trốn thoát khỏi khu ổ chuột vì có nhiều người sẵn sàng giúp ông.\\[4pt]
\textit{(Janusz Korczak was a famous Jewish doctor and educator in Warsaw; he could have escaped the ghetto as many were willing to help him.)}

Nhưng khi quân Đức đến giải các em nhỏ trong trại trẻ mồ côi của ông đi Treblinka, ông từ chối rời đi.\\[4pt]
\textit{(But when German soldiers came to take the orphan children from his care to Treblinka, he refused to leave.)}

Ông dẫn 192 trẻ mồ côi tiến ra sân tập hợp --- mặc quần áo đẹp nhất, mỗi tay cầm một bông hoa --- đi đến chỗ chết với đầu ngẩng cao.\\[4pt]
\textit{(He led 192 orphan children into the assembly yard --- dressed in their best clothes, each holding a flower --- walking to their deaths with heads held high.)}

Một người lính Đức hỏi: ``Ông là bác sĩ Korczak phải không? Ông có thể đi.'' Korczak nhìn những đứa trẻ rồi lắc đầu.\\[4pt]
\textit{(A German soldier asked: `Are you Dr. Korczak? You may go.' Korczak looked at the children and shook his head.)}

\end{truyen}
\begin{baihoc}
  Tình yêu thương thật sự đặt người ta yêu lên trên cả sinh mạng của mình --- đó là đỉnh cao của phẩm giá con người.\\[4pt]
  \textit{(True love places those we love above our own life --- that is the pinnacle of human dignity.)}
\end{baihoc}

\section{Người Mẹ Sói Và Bài Học Buông Tay}
\begin{truyen}{Người Mẹ Sói Và Bài Học Buông Tay}{Thiên nhiên --- Bài học về Tình Mẫu Tử}
\chuhoa{S}{ói} mẹ dạy sói con săn mồi bằng cách để chúng vấp ngã và thất bại --- không bao giờ nhảy vào cứu khi con đang vật lộn với con mồi.\\[4pt]
\textit{(A mother wolf teaches her cubs to hunt by letting them fall and fail --- never leaping in to help when they are struggling with prey.)}

Khi cần, sói mẹ sẽ giảng giải, hướng dẫn, làm mẫu --- nhưng sẽ không làm thay.\\[4pt]
\textit{(When needed, the mother wolf will explain, guide, and demonstrate --- but will never do it for them.)}

Những con sói non được mẹ làm hết thay sẽ không sống được khi một mình.\\[4pt]
\textit{(Wolf cubs whose mother does everything for them will not survive when alone.)}

Yêu con không phải là bảo vệ con khỏi mọi khó khăn --- mà là chuẩn bị cho con đủ sức đối mặt với mọi khó khăn.\\[4pt]
\textit{(Loving your child is not protecting them from all hardship --- it is preparing them with enough strength to face all hardship.)}

\end{truyen}
\begin{baihoc}
  Tình yêu thương thật sự đôi khi có nghĩa là đứng nhìn và để người ta tự đứng dậy --- vì đó là cách họ trưởng thành.\\[4pt]
  \textit{(True love sometimes means standing by and letting someone stand up on their own --- for that is how they grow.)}
\end{baihoc}
